\clearpage\begingroup\SourceLeaf{51}\NoSourceHead
\begin{center}
[D.]\\[.7em]
\textbf{[VERMISCHTES.]}\\[1.2em]
[1.]\\[.7em]
\textsc{Gauss} an \textsc{Bessel}. Göttingen, 27. Januar 1829.
\par\smallskip\rule{.72\textwidth}{.35pt}\\
{\small [Briefwechsel zwischen \textsc{Gauss} und \textsc{Bessel}, Berlin 1880, S. 488.]}\\
\rule{.72\textwidth}{.35pt}
\end{center}
\noindent .\ .\ .\ .\ .\ .\ .\ .\ .\ .\ .\ .\ .\ .\ .\ .\ .\ .\ .\ .\ .
\par\noindent\hspace*{\parindent}Erst ganz seit Kurzem habe ich endlich einmal wieder anfangen können,
an eine wissenschaftliche Arbeit zu denken. Ich habe verschiedene schon vor
langer Zeit gehabte Ideen wieder aufgenommen, die sich auf eine von der
\textsc{Laplace}schen ganz verschiedene Behandlung der Theorie der Capillaraction
beziehen, und es ist mir bereits gelungen, das Meiste von dem, was mir früher
nur als möglich vorschwebte, zu verwirklichen. Ich hoffe auch mit dem noch
fehlenden Theile schon fertig zu werden, und wenn ich nur nicht zu früh
wieder anderweitige Störungen erfahre, diesen Gegenstand schriftlich auszuführen,
und schmeichle mir, dass Sie die Ausbeute nicht ohne Interesse demnächst
lesen werden.

Bei dieser Veranlassung habe ich nun aber auch die beiden Abhandlungen
von \textsc{Laplace} \textit{Th. de l’act. cap.}[\textsuperscript{*})]\begingroup\renewcommand{\thefootnote}{}\footnotetext{[\textsuperscript{*}) P. S. \textsc{Laplace}, \textit{Théorie de l’action capillaire}, Supplément au livre dixième de la \textit{Mécanique céleste}, Paris 1806; Oeuvres IV, 1880, S. 349.]}\endgroup\ und \textit{Suppl. à la th. de l’act. cap.}[\textsuperscript{**})]\begingroup\renewcommand{\thefootnote}{}\footnotetext{[\textsuperscript{**}) P. S. \textsc{Laplace}, \textit{Supplément à la théorie de l’action capillaire}, Paris 1807; Oeuvres IV, 1880, S. 419.]}\endgroup\ 
sorgfältiger studirt, als früher von mir geschehen war. Dabei bin ich aber
zu meiner Verwunderung auf etwas gestossen, worüber ich gern Ihre Ansicht
erfahren möchte, wenn Sie anders diesen Gegenstand selbst auch zu einem
Gegenstand Ihres Studiums gemacht haben.
\clearpage\SourceLeaf{52}\HistoricalHead{47}{R}{GAUSS AN BESSEL.}
\noindent\hspace*{\parindent}Im Grunde hat \textsc{Laplaces} Theorie zwei Cardinalpuncte:
\begin{list}{}{\setlength{\leftmargin}{2.5em}\setlength{\labelwidth}{1.5em}\setlength{\labelsep}{.5em}\setlength{\itemsep}{0pt}\setlength{\topsep}{0pt}}
\item[1)] die Wirkung eines Meniskus auf eine unendlich enge Röhre, welche
er auf eine sehr schöne Art entwickelt hat, und woraus leicht die
Fundamentalgleichung für die Oberfläche folgt;
\end{list}
\noindent\begin{minipage}[t]{.80\textwidth}\setlength{\parindent}{1.2em}
\begin{list}{}{\setlength{\leftmargin}{2.5em}\setlength{\labelwidth}{1.5em}\setlength{\labelsep}{.5em}\setlength{\itemsep}{0pt}\setlength{\topsep}{0pt}}
\item[2)] den Satz, dass der Winkel der Oberfläche der Flüssigkeit
mit der Gefäss-Wand constant ist (bei einerlei Flüssigkeit
und einerlei Gefäss-Materie).
\end{list}
\noindent\hspace*{\parindent}Dieser zweite Satz wird überall in der Abhandlung vorausgesetzt.
\end{minipage}\hfill
\begin{minipage}[t]{.17\textwidth}\vspace{0pt}\centering
\begin{tikzpicture}[x=1cm,y=1cm,line width=.55pt]
\draw (0,2.1) .. controls (.4,1.99) and (.8,1.99) .. (1.2,2.1) -- (1.2,1.94) -- (0,1.94) -- cycle;
\draw (.60,1.94) -- (.60,0);
\end{tikzpicture}
\end{minipage}
\par\smallskip
\begin{list}{}{\setlength{\leftmargin}{5.2em}\setlength{\labelwidth}{4em}\setlength{\labelsep}{.5em}\setlength{\itemsep}{.3em}\setlength{\topsep}{0pt}}
\item[pag. 20] »inclinaison qui, \textls[120]{comme on l’a vu}, doit être la même pour
tous les plans«
\item[pag. 23] »et \(\theta'\) est, \textls[120]{comme on l’a vu}, une quantité indépendante de ce
demi-diamètre«
\item[pag. 14] in der zweiten Abh[andlung] »l’angle \(\tilde\omega\) est constant, comme je
\textls[120]{l’ai fait voir} dans la théorie citée«.
\end{list}
\noindent Allein in der ganzen ersten Abhandlung \textls[120]{selbst} finde ich kein Wort, was
dienen kann dies \textls[120]{zu beweisen}. Es kann also wol nichts gemeint sein als
die Stelle in der Einleitung pag. 5, wo ich aber den Schluss, dass die »plans
(en question) sont également inclinés à leur parois« keinesweges auf eine befriedigende
Art begründet finde. Ich gestehe, dass mir dieser \textls[120]{Haupttheil}
von \textsc{Laplaces} Theorie der präcisen mathematischen Begründung des übrigen
keinesweges würdig zur Seite zu stehen, sondern mehr den Charakter der vagen
Aperçus, die man früher von dem ganzen Phänomene hatte, zu tragen scheint.

Freilich könnte man sagen, dass \textsc{Laplace} diese Lücke einigermassen in
der zweiten Abhandlung ausgefüllt hat. Das Rapprochement der ersten Methode
die Haarröhrchen zu behandeln mit der anderen in der zweiten Abhandlung
(die doch wohl im Grunde nichts weiter ist als die \textsc{Lalande}sche)
führt zu einer Bestimmung des Winkels quaestionis pag. 18. Allein erstlich
ist dies doch wol ein unwürdiger Umweg, und zweitens folgt\textsuperscript{*})\begingroup\renewcommand{\thefootnote}{}\footnotetext{\textsuperscript{*}) Auch möchte sich gegen diesen Beweis wol sonst noch allerlei erinnern lassen.}\endgroup\ das Constantsein
des Winkels daraus nur für den ganz speciellen Fall, dass die Gefässwände
einen \textls[120]{senkrechten} Cylinder bilden.

Sie würden mich sehr verpflichten, wenn Sie mir Ihre Ansicht hierüber
gefälligst bald mittheilen wollten. Es ist aber nicht die Rede davon, in wie
\HistoricalHead{47}{R}{GAUSS AN BESSEL.}\clearpage\SourceLeaf{53}\HistoricalHead{48}{L}{PHYSIK. BRIEFWECHSEL.}
\noindent fern man dies vage Räsonnement pag. 5 der Einleitung durch tiefere Untersuchung
zu der Würde eines Beweises erheben kann (denn das ist meine eigene
Meinung sogar), sondern ob nach Ihrem Gefühl jenes \textls[120]{wie es da steht} die
gerechten Forderungen, die der Geometer macht, befriedigt. Die Untersuchungen
pag. 44 sqq. geben eigentlich für den beabsichtigten Zweck gar
kein Resultat.\quad .\ .\ .\ .\ .\ .\ .\ .\ .\ .\ .\ .\ .\ .\ .\ .\ .\ .\ .\ .\ .
\clearpage\endgroup
